\chapter{Il sistema}
\label{cap:sistema}

Questo capitolo descrive l'architettura hardware dell'esacottero e ne riepiloga
la composizione tramite la \emph{Bill of Materials} (BOM). L'obiettivo \`e
fornire il quadro fisico del velivolo su cui poggiano tutti i capitoli
successivi: la configurazione del firmware, l'acquisizione dati e la diagnosi
dei guasti si riferiscono ai sottosistemi qui introdotti.

\section{Architettura generale}

Il velivolo \`e un \textbf{esacottero}, ossia un multirotore a sei rotori
disposti simmetricamente sui bracci del frame. Rispetto a un quadricottero,
la configurazione a sei motori offre maggiore capacit\`a di carico e una
ridondanza intrinseca: la perdita parziale di spinta su un rotore pu\`o essere
in parte compensata dagli altri. Questa caratteristica \`e rilevante per il
progetto, perch\'e la campagna di volo (Capitolo~\ref{cap:campagna}) studia
proprio l'effetto di una pala danneggiata sul comportamento del velivolo.

I sottosistemi principali sono cinque:
\begin{description}
  \item[Struttura] il frame F550 con carrello e mast GPS;
  \item[Propulsione] sei gruppi motore-ESC-elica;
  \item[Avionica] il flight controller Pixhawk 6X con IMU, magnetometro e GNSS;
  \item[Alimentazione] batteria LiPo, distribuzione di potenza e regolatori;
  \item[Comunicazione] radiocomando RC e telemetria verso la ground station.
\end{description}

\figurasegnaposto{Schema a blocchi dell'architettura: Pixhawk 6X come nodo
centrale collegato a GPS (UART), sei ESC con telemetria DShot, ricevitore RC,
telemetria 433\,MHz e power module.}{Schema -- da generare}{fig:schema-blocchi}

\figurasegnaposto{Vista del drone con i componenti principali etichettati.}{Foto
-- da scattare}{fig:drone-etichettato}

\section{Sottosistemi}

\subsection{Struttura e propulsione}
Il frame DJI F550 \`e un telaio esagonale che integra la \emph{Power
Distribution Board} (PDB) nella piastra inferiore, semplificando il cablaggio
di potenza verso i sei ESC. Il carrello \`e costituito da piedi flessibili ad
arco che ammortizzano l'atterraggio. La propulsione \`e affidata a sei motori
brushless \textbf{AIR2213 KV920} (tre in rotazione oraria e tre antioraria,
necessari per l'equilibrio di coppia del multirotore), pilotati da sei ESC
\textbf{Tekko32 F4} con telemetria DShot. Quest'ultima \`e particolarmente
utile ai fini del progetto, perch\'e permette di registrare regime, corrente e
temperatura di ciascun motore (Capitolo~\ref{cap:acquisizione}).

\subsection{Avionica e navigazione}
Il flight controller \`e un \textbf{Pixhawk 6X}, basato su microcontrollore
STM32H7 e progettato per il firmware PX4. Integra \textbf{tre IMU ridondanti}
(accelerometro + giroscopio) gestite in \emph{sensor voting}: il firmware
confronta le letture e seleziona/scarta i sensori in caso di discrepanza,
aumentando l'affidabilit\`a della stima d'assetto. La navigazione si appoggia
a un modulo \textbf{GNSS CubePilot Here+} (con magnetometro integrato) montato
su un mast di circa 15\,cm per allontanarlo dalle interferenze
elettromagnetiche generate dai cavi di potenza e dai motori. Il sistema
supporta la correzione \textbf{RTK} tramite una stazione base, per una
precisione di posizionamento centimetrica.

\subsection{Alimentazione}
L'energia \`e fornita da una batteria \textbf{LiPo 4S da 5600\,mAh}
(16{,}8\,V a piena carica, connettore XT60). La PDB distribuisce la potenza ai
sei ESC; un regolatore step-down (BEC) ricava i 5\,V per l'avionica e le
periferiche. Il monitoraggio di tensione e corrente avviene tramite un power
module collegato al Pixhawk, la cui corretta configurazione \`e stata oggetto
di un intervento specifico (Capitolo~\ref{cap:configurazione}).

\subsection{Comunicazione e controllo}
Il controllo manuale avviene tramite un radiocomando FrSky con ricevitore
\textbf{X8R} a bordo. Il collegamento dati con la ground station
\textbf{QGroundControl} \`e realizzato da una coppia di moduli telemetria
\textbf{Holybro a 433\,MHz}, che trasmette telemetria in tempo reale e consente
la modifica dei parametri e il monitoraggio dei failsafe.

\section{Sintesi della Bill of Materials}

La Tabella~\ref{tab:bom} riepiloga i componenti principali con i codici
identificativi adottati nel repository (S struttura, P propulsione, A avionica,
E alimentazione, C comunicazione, SW software). Il dettaglio completo, con note
e voci ancora in via di definizione, \`e mantenuto nel file \texttt{docs/BOM.md}.

\begin{table}[H]
\centering
\caption{Estratto della Bill of Materials dell'esacottero F550.}
\label{tab:bom}
\small
\begin{tabularx}{\linewidth}{@{}l L{3.2cm} L{3.6cm} c@{}}
\toprule
\textbf{ID} & \textbf{Componente} & \textbf{Modello} & \textbf{Q.t\`a} \\
\midrule
S-01 & Frame & DJI F550 (PDB integrata) & 1 \\
S-02 & Carrello & Piedi flessibili ad arco & 1 set \\
S-03 & Mast GPS & Albero $\sim$15\,cm anti-EMI & 1 \\
\addlinespace
P-01 & Motore brushless & AIR2213 / KV920 & 6 \\
P-02 & ESC & Tekko32 F4 (telemetria DShot) & 6 \\
P-03 & Eliche & 2 pale (dim.\ da confermare) & 6 \\
\addlinespace
A-01 & Flight controller & Pixhawk 6X (PX4, 3 IMU) & 1 \\
A-06 & GPS / GNSS & CubePilot Here+ (su mast) & 1 \\
A-07 & Base RTK & Stazione base & 1 \\
A-09 & Cavo GPS$\leftrightarrow$Pixhawk & Auto-costruito, schermato (JST-GH) & 1 \\
\addlinespace
E-01 & Batteria & LiPo 4S 5600\,mAh (XT60) & 1 \\
E-02 & BEC / regolatore & Step-down 16{,}8\,V $\rightarrow$ 5\,V & 1 \\
\addlinespace
C-01 & Ricevitore RC & FrSky X8R & 1 \\
C-03 & Telemetria & Holybro 433\,MHz (TX/RX) & 1 set \\
\addlinespace
SW-01 & Firmware FC & PX4 & --- \\
SW-02 & Ground station & QGroundControl & --- \\
\bottomrule
\end{tabularx}
\end{table}

\begin{notabox}[Nota sulla configurazione hardware]
Il flight controller attuale (Pixhawk 6X) \`e il risultato di una sostituzione
hardware resa necessaria da un guasto del precedente Cube Black
(Capitolo~\ref{cap:diagnostici}). Anche il cavo GPS \`e stato ricostruito con
schermatura dedicata a seguito della diagnosi dei dropout
\texttt{sensor\_gps}. La BOM riflette quindi lo stato hardware consolidato a
valle di questi interventi.
\end{notabox}
